\section{Introduction}
\label{sec:introduction}
The capacity to cross interstellar distances does not settle whether a
civilization will keep paying to do so. Yet an inward explanation of the
Fermi paradox faces a stronger objection than this elementary distinction:
resource growth itself can renew the incentive to expand. A civilization
that values additional computation need not exhaust its ambitions before
acquiring another star. If extensive industry produces resources and
independent options that intensive local investment cannot supply cheaply,
better computers may strengthen neither local sufficiency nor the decision
to remain local. The computational inward transition hypothesis (\CITH)
therefore needs a boundary that expansion can cross, not a verbal claim
that virtual experience is attractive.

Expansion arguments associated with \citet{Hart1975} and \citet{Tipler1980}
connect technological capability to persistent settlement or self-replicating
activity. The wider assumptions behind the Great Silence were surveyed by
\citet{Brin1983}. Settlement dynamics, finite site lifetimes, and spatial
heterogeneity complicate complete occupation \citep{NewmanSagan1981,CarrollNellenback2019}.
They do not make launch behaviour irrelevant: an expansionist lineage can
have consequences far beyond its original population. Our question is
whether an internal computational alternative removes the incentive to
create that lineage, and whether removing it in many civilizations is enough
to explain no arrival in a specified region.

The answer developed here begins against the strong hypothesis. With
logarithmic computational value, concave local production, bounded external
prices, and a positive residual external motive, unbounded physical budgets
eventually exclude all-inward allocation, whatever local computational
productivity does. This result treats efficiency and physical resources as
different quantities. A processor improvement changes the useful service
obtained from a fixed energy or material allocation; access to another
resource base changes the allocation that can be made. Conflating these
changes would remove precisely the outward pressure that the model tests.

\subsection{Contribution and organizing argument}
The contribution has four connected parts. First, a common corner criterion
yields a tight production-independent floor and a unique tipping budget
under logarithmic value. An explicit resource-acquisition process has a
unique initial-productivity boundary separating eventual inward investment
from outward lock-in, with a critical path approaching the tipping budget.
Second, optimizing non-replicating observation and dormant backups changes
the relevant corner and replaces gross project prices by opportunity costs
net of substitute services. A growing-budget construction shows that the two corners can even have opposite persistence conclusions. A catastrophe model then derives the resilience
premium from value at risk rather than assigning it freely. Third, a
free-world-number margin connects a specified population ethic to value
curvature, while an infinite-horizon benchmark and subjective-delay
inequalities make continuation returns explicit. Finally, a marked-origin
arrival model quantifies how small the expansionist tail must be when
potential arrival exposure is large.

These results organize a particular inward scenario around three tests:
a viable resource path that stays on the appropriate side of an incentive
boundary, credible non-replicating ways of satisfying external information
and resilience motives, and sufficiently weak time-adjusted returns on
replication. The tests must hold together in a scenario relying on all
three mechanisms. They are not universal necessary conditions for every
possible inward civilization. Slower marginal-value decay, vanishing
external preferences, strongly increasing local returns, or rising external
prices supply different routes. The model proves conditional exclusions
and existence results, not a universal evolutionary destiny. Its strongest
negative conclusion is correspondingly precise: within the logarithmic,
concave-production benchmark, productivity alone cannot defeat a persistent
residual motive when the physical budget grows without bound.

The corner calculation uses standard concave optimization. Its contribution
is not a new first-order optimality theorem. The additional theory lies in
sharp feasibility across production functions, the pathwise obstruction,
the unique basin boundary, the exact scope and failure of an exponential
endpoint test, and the links between service substitution, intertemporal
returns, and population exposure. Each link has its own assumptions. In
particular, a sufficient condition for one civilization to stop new
investment is not a theorem that all existing autonomous descendants stop
replicating. A population argument must keep their exposure in the record.

\subsection{What is inherited from prior work}
Computational and postbiological civilizations are established subjects,
not ideas introduced by this paper \citep{Dick2003,CirkovicBradbury2006}.
\citet{Smart2012} describes movement toward dense internal scales. Our
increasing-returns result supplies a resource-response condition against
which that kind of mechanism can be tested. It does not import the stronger
claim that advanced civilizations invariably leave the observable universe.
Bradbury's Matrioshka-brain manuscript describes stellar computation and
nested heat-use architectures \citep{Bradbury1999}; the thermal calculation
below is a deliberately smaller bridge between allocation and search,
not a proposal of that architecture.

Aestivation instead concerns when to compute \citep{Sandberg2017}. The
thermodynamic objection of \citet{BennettHansonRiedel2019} emphasizes that
available entropy reservoirs must be specified before efficiency advantages
are inferred. Our latency comparison concerns delayed external returns
relative to ongoing internal activity, and does not assume that waiting
for a colder universe is optimal. Sustainability and homeostatic proposals
address other reasons expansion might end \citep{HaqqMisraBaum2009,WongBartlett2022}.
They can coexist with the allocation mechanism but do not prove its
thresholds. Conversely, self-replicating resource acquisition
\citep{ArmstrongSandberg2013} and population-total welfare
\citep{Bostrom2003Waste} are direct counterpressures. They are included as
cases in which outward investment can remain attractive.

\subsection{Three meanings of inwardness}
We reserve \emph{all-inward allocation} for zero new expenditure on resident,
industrial, and observation projects. \emph{No replication} instead permits
observation and dormant backups, while setting new reproducing capacity
to zero. \emph{Low detectability} is an observational property and can hold
or fail under either allocation. These definitions also distinguish a
civilization's current investment from the continued activity of earlier
launches. A society may stop launching because its distant industry is
already reproducing effectively; that is not an explanation of silence.

The paper proceeds through these distinctions rather than treating every
external project as colonization. The model first separates physical
support from economic preference, then identifies the corner and its
persistence conditions. Residual motives and latency determine whether
those conditions survive a more explicit account of services and time.
The population calculation finally asks how much outward or leaked activity
a chosen null likelihood tolerates. Detailed workload examples, the paired
sensitivity ensemble, and the original two-period construction are in the
supplement. Their numerical values are preserved as implementation checks,
not used as a prior over unknown civilizations.
